\chapter{Discussion and Extensions}
\label{chap:discussion}

\section{What the Implemented Connection Establishes}
The principal result is an implemented connection between a human-specified local task and an executable motion--force skill. The teaching representation supplies pose, loading, and timing together. The policy consumes observations and generates that same kind of sequence. The common executor and existing contact controller make the output actionable on the robot. E1 examines strategy execution, E2 separates target-force learning and force conditioning, and E3 measures image-based ink-area reduction after spatial reuse.

The primary E1 comparison establishes that this connection produces an executable local contact skill. Rules, replay, and the learned provider each complete contact, the required loaded traversal, and exit in all five planned runs. Their execution characteristics differ: T has a shorter mean completion time and smaller time dispersion than C, and C has greater mean excess-force exposure. E2 then shows a shorter operator-based completion interval for the force-observing configuration than for B, together with an incomplete release transition in A4. E3 adds condition-mean ink-area reductions of \EthreeDtwoRemoval--\EthreeNRemoval\% in the evaluable registered placements. Together, the evaluations support an executable skill, observed force-configuration differences, and spatial reuse with its explicit adapter, while giving fixed strategies their observed value.

The comparison also clarifies what an execution-level result can establish. A generated rise can be learned as part of the demonstrated temporal structure. A gate can recognize unloading, and an evaluator can confirm a displacement. None of those observations directly measures remaining defect depth. Maintaining that distinction allows the method to be described positively and precisely without assigning hidden process reasoning to the policy.

\section{When Rules and Replay Are Sufficient}
The nominal task is narrow: a specified region, a single demonstrated direction, a fixed setup, and an existing lower-level controller. Under those conditions, a prescribed strategy has a strong opportunity to work well. T reuses a full force-inclusive episode and its timing, making it a substantive alternative rather than a deliberately impoverished baseline. Its short completed intervals and low observed dispersion indicate that a stable demonstration can be valuable without policy learning.

R illustrates the same point from a process-recipe perspective. A finite rule can express approach, loading, traversal, and departure directly. Its longer interval reflects the constructed strategy and contact synchronization, not an intrinsic limitation that all rules must be slower. The executed additional \(R^{*}\) diagnostic makes this sensitivity concrete: processing-only retiming is associated with a \RstarTimeGain~s (\RstarTimePercent\%) shorter mean interval than original R. The dataset mean was used as a quintic peak-feed parameter, not matched actual average speed. The separate date and one-then-four collection history limit causal interpretation.

Imitation learning is used here to represent the demonstrated temporal relationship between motion and target force as an observation-conditioned generator. The specified region alone does not determine that relationship, but rules and replay can also supply it. E1 consequently evaluates learning alongside those alternatives rather than treating it as a necessary condition for completing the fixed task. T's strong nominal baseline and C's executable sequence are compatible findings: they characterize different ways to supply the local strategy.

\section{Roles of Human Guidance and Relative Representation}
Human guidance resolves where the local operation is intended to occur. Teaching supplies an example of how movement and loading can be organized. These are complementary contributions of the operator. A spatial mask, mark, or origin cannot be treated as a complete polishing strategy, and a demonstrated strategy cannot silently substitute for a new region specification.

The relative representation supports this division by learning position around a stain center and explicitly handling the stain direction at the policy boundary. A single-direction local movement pattern is reused by aligning live XY and the policy image with the teaching direction, then inversely rotating the generated XY path and adding the live center. This is an implemented spatial-generalization mechanism, rather than a future proposal. It addresses placement and in-plane direction through task registration while leaving observation-conditioned motion--force generation to the learned policy.

Its scope follows from the transformed channels. Z, tool orientation, and force are not rotated, and image adaptation is an in-plane warp about the interaction ROI. Thus the mechanism transfers the planar movement pattern without claiming complete pose--wrench equivariance or arbitrary-view invariance. E1/E2 disable direction alignment; E3 uses it for D1/D2 and evaluates the registered positions/directions in Table~\ref{tab:e3_placements}. Those observations characterize the combined policy and adapter in the tested planar arrangement.

The fixed image interaction ROI also has a specific function. It focuses part of the visual representation around the calibrated tool interaction area while retaining a global feature. It is not an explicit defect-state channel. The current experiments evaluate the resulting observation path as a whole; they do not isolate whether a mark's appearance, pose, or force carries most of the policy's decisions.

\section{What Force Output and Observation Contribute}
Target-force learning allows the reference sequence to express a changing load, including a departure from a constant recipe. Measured-force observation provides information about the interaction that occurred rather than the load merely requested. Their roles are conceptually distinct and are separated by the E2 contracts.

The E2 observations identify differences in the organization of contact and release. C's mean operator-based completed interval is \BCSeconds~s (\BCPercent\%) shorter than B's, with five completed attempts in each. This is a descriptive comparison of separately trained checkpoints, not the same weights tested with a channel switched off. A4 exposes an incomplete transition in the motion-only condition, but one interrupted attempt does not identify missing force output as the cause. The generated motion, external constant-force interface, mechanical state, and timing can all contribute.

The magnitude of force peaks also limits an uncomplicated advantage claim. C does not reduce the E2 mean peak below A, and it has greater E1 excess exposure than R and T. A lower whole-window force mean can follow from less time spent in contact. It is therefore inadequate to call that lower mean an improvement in force-control accuracy or surface finishing.

The force-frame convention is another practical boundary. The resolved teaching-to-base relation supports empirical common-axis analysis, while the evaluated training file and online adapter retain their actual numerical conventions. The force-conditioning path was not redefined for this thesis. Its present performance is an observation about that implementation; a future frame-consistent learning comparison would be a separate experiment. Similarly, the executor suppresses tangential force commands, so a nine-dimensional learned action should not be described as a physically demonstrated three-axis force-tracking capability.

\section{Procedural Completion and Surface Quality}
The contact procedure is a meaningful intermediate target for local polishing. A tool must establish contact, move through the relevant region under load, and leave the surface. E1 provides explicit observations for that sequence. E2 contributes an operator-based interpretation of task progression that preserves the interrupted attempt.

E3 adds a directly observed processing outcome: the percentage reduction of segmented board-marker ink area. N/P1/D1/D2 have means of \EthreeNRemoval, \EthreePoneRemoval, \EthreeDoneRemoval, and \EthreeDtwoRemoval\% across \EthreeImageN{} evaluable image pairs. These high condition means support practical reuse of the same policy with task registration and direction alignment within the tested placements. The 83.879\% D2-3 value remains an individual observation, while D1-4 and P2 remain image-unavailable; the range of means is not a minimum performance guarantee.

A completed traversal, operator completion, and a reduced image stain are complementary observations. E3 provides a quantitative ink-area outcome for its own spatial-transfer cohort. It cannot establish E1/E2 removal quality or productivity at equal removal. Neither marker disappearance nor an execution-derived heatmap independently measures metal depth, roughness, or micro-scratch restoration. This distinction preserves the positive result while stating its physical scope.

\section{Development Extensions}
Execution-derived dose heatmaps, finite blending paths, and adaptive reference-clock execution are implemented development assets. The heatmap remains an uncalibrated diagnostic; the blending and adaptive records do not provide a controlled E1--E3 surface-quality comparison. High-level residual reinforcement learning for curved-surface restoration remains planned. Appendix~\ref{app:development_extensions} preserves the diagnostic formulation, example heatmap, version distinctions, and implementation-status table. These extensions are separate from the three evaluated contributions.

\section{Scope and Limitations}
E1/E2 assess the nominal local task; E3 extends the evidence to the reported planar placements and directions. Arbitrary shapes, curved surfaces, height changes, and reference-Z robustness remain outside those observations. The tested direction transfer includes an explicit observation/image adapter and XY inverse transformation. It is not an isolated test of policy equivariance or of force observation's causal contribution to spatial transfer. The inherited blending/RL assets retain their separate implementation or planned status.

There is one recorded training seed per learned condition, distinct from the repeated robot executions. Original E1 shuffles R/T/C within five blocks, whereas \(R^{*}\) is a separate-date one-then-four diagnostic. E2 runs A, then B, then C, using separate checkpoints; A/B also changes action dimension and force-target composition. Small cohorts and sample SDs establish neither statistical significance nor equivalence or population reliability. Setup, execution order, and operator timing remain possible contributors to the observed differences.

E3 is a selected development collection with retained retries and dates. N uses the existing HOME return, P1 uses PTP, and D1/D2 add commanded-release confirmation to PTP. These differences preclude treating duration or path length as a controlled removal-efficiency ranking. The image analysis is post hoc and semiautomatic, using stored photographs and saved visible-boundary annotations. D1's removal mean has \(n=4\); all five P2 final photographs and D1-4 are occluded. Hidden stain area is not estimated. The common 40/45/50 threshold sensitivity check describes segmentation dependence, not an uncertainty interval for future robot performance.

The numerical evidence is traceable to frozen run identities and archived analysis inputs. E1/E2 precision aggregates agree with their source CSVs; the additional \(R^{*}\) retains unrounded per-run values and template/configuration identifiers. E3 reprocessing from the original RGB, saved annotations, and fixed ROI reproduces all 25 initial areas, 19 evaluable final areas, and 75 stored mask/ROI arrays. This verifies computational reproduction of the established image analysis. It does not independently validate its manual semantic labels or calibrate physical material removal.

E1 automatic completion, E2 operator completion, and E3 image-area reduction remain distinct. A4's successful time is null; its \AfourInterval~s interruption and all-attempt force/path values remain. C5 stays in the original C group; C6 is additional. Metal depth, mass, roughness, and micro-scratch restoration are not independently measured, and an unrecorded controller-applied target still prevents force-tracking-error claims.
